\chapter{Laxatives}
\index{Laxatives}

\section*{Definition and Overview}
Laxatives are substances or medications that loosen stools or stimulate bowel movements to treat or prevent constipation. They are categorized based on their mechanism of action, including bulk-forming, osmotic, stimulant, and stool softener (emollient) laxatives. While widely available over-the-counter and generally safe for short-term use, their selection must be tailored to the individual's clinical circumstances, as inappropriate use or overuse can lead to serious electrolyte imbalances and chronic bowel dysfunction.

\section*{Epidemiology}
Constipation is a highly prevalent gastrointestinal complaint, affecting approximately 15\% to 20\% of the global population and a similar proportion in many countries. It is more common in women, older adults (particularly those aged 65 years and over), and pregnant individuals. Consequently, laxatives are among the most frequently purchased over-the-counter and prescribed medications. Laxative overuse or misuse is particularly common in elderly populations due to age-related bowel changes, as well as in individuals with eating disorders who incorrectly use stimulant laxatives for weight management.

\section*{Aetiology and Pathophysiology}
Laxatives act through distinct physiological mechanisms to increase stool water content, stimulate colonic motility, or facilitate stool passage.
\begin{itemize}
    \item \textbf{Bulk-forming laxatives:} These agents, such as psyllium or methylcellulose, absorb water in the intestinal lumen, increasing stool mass and softening consistency to naturally stimulate peristalsis.
    \item \textbf{Osmotic laxatives:} Compounds like lactulose, macrogol (polyethylene glycol), and magnesium salts draw water into the bowel through osmotic gradients, increasing stool volume and softening it for easier passage.
    \item \textbf{Stimulant laxatives:} Agents such as senna and bisacodyl directly irritate the mucosal lining or stimulate the myenteric plexus, increasing intestinal motility and active fluid secretion.
    \item \textbf{Stool softeners (surfactants):} Emollients like docusate sodium lower the surface tension of the stool, allowing water and lipids to penetrate and soften the faecal mass.
    \item \textbf{Lubricant laxatives:} Liquid paraffin coats the stool and intestinal wall, reducing water absorption and easing the passage of hard stools through the rectum.
\end{itemize}

\section*{Clinical Features}
The administration of laxatives typically results in predictable changes in bowel habits, but inappropriate use can present with systemic features.
\begin{itemize}
    \item \textbf{Therapeutic effects:} Increased frequency of bowel movements, softer stool consistency, and reduced straining during defaecation.
    \item \textbf{Gastrointestinal symptoms:} Mild abdominal cramping, bloating, flatulence, and abdominal distension, particularly common with bulk-forming or osmotic agents.
    \item \textbf{Stimulant-induced cramping:} Sharp, colicky abdominal pain occurring shortly after ingestion, characteristic of stimulant laxatives.
    \item \textbf{Dehydration and electrolyte disturbance:} Lethargy, muscle weakness, palpitations, and dizziness, which occur when laxative overuse leads to excessive watery diarrhoea.
\end{itemize}

\section*{Differential Diagnosis}
Before initiating laxative therapy, it is crucial to differentiate simple functional constipation from other underlying medical conditions.
\begin{itemize}
    \item \textbf{Mechanical obstruction:} Colorectal cancer, strictures, or volvulus can block the bowel, presenting with constipation but requiring surgical intervention rather than laxatives.
    \item \textbf{Endocrine disorders:} Hypothyroidism and diabetes mellitus can significantly slow bowel motility.
    \item \textbf{Neurological disorders:} Parkinson's disease, multiple sclerosis, and spinal cord injuries can impair defaecation reflexes.
    \item \textbf{Medication-induced constipation:} Opioid analgesics, calcium channel blockers, iron supplements, and anticholinergic agents are common causes of secondary constipation.
    \item \textbf{Irritable bowel syndrome (IBS-C):} A functional disorder characterized by constipation accompanied by chronic abdominal pain and bloating, requiring a comprehensive management plan beyond simple laxatives.
\end{itemize}

\section*{When to Seek Medical Care}
Patients should seek medical advice if they experience persistent constipation that does not respond to dietary adjustments or short-term laxative use, or if they require laxatives for more than a week.

\textbf{Contact your local emergency services for:}
\begin{itemize}
    \item Severe, sudden, or worsening abdominal pain
    \item Blood in the stool or rectal bleeding
    \item Unexplained weight loss combined with changes in bowel habits
    \item Signs of severe electrolyte imbalance, such as confusion, palpitations, or fainting
\end{itemize}

\section*{Diagnosis and Investigations}
The diagnosis of constipation is primarily clinical, based on patient history and criteria such as the Rome IV criteria, but investigations are warranted if red flag symptoms are present or if constipation is chronic and refractory.
\begin{itemize}
    \item \textbf{Clinical history and examination:} Detailed assessment of diet, fluid intake, medication use, physical activity, and a digital rectal examination to check for faecal impaction.
    \item \textbf{Blood tests:} Thyroid function tests to rule out hypothyroidism, serum calcium to rule out hypercalcaemia, and a full blood count to screen for anaemia.
    \item \textbf{Abdominal X-ray:} Helpful in acute cases to assess the degree of faecal loading or to check for bowel obstruction.
    \item \textbf{Colonoscopy:} Indicated in patients over the age of 50, or those with alarm features (e.g. rectal bleeding, weight loss, family history of bowel cancer), to exclude malignancy.
    \item \textbf{Anorectal manometry and transit studies:} Specialized investigations to assess sphincter function and colonic transit time in refractory cases.
\end{itemize}

\section*{Management}
The management of constipation should focus on lifestyle modifications first, with laxatives used as a supportive, step-wise pharmacological measure under clinical guidance.
\begin{itemize}
    \item \textbf{Lifestyle modifications:} Increasing dietary fibre intake (to 25-30g daily), maintaining adequate hydration, and engaging in regular physical exercise.
    \item \textbf{First-line pharmacological therapy:} Bulk-forming laxatives are generally preferred initially as they mimic natural physiological processes.
    \item \textbf{Second-line osmotic agents:} Macrogol or lactulose are introduced if bulk-forming agents are ineffective or poorly tolerated.
    \item \textbf{Short-term stimulant use:} Senna or bisacodyl may be added for acute relief or in opioid-induced constipation, but should be used sparingly.
    \item \textbf{Rectal therapies:} Glycerol suppositories or microenemas can be utilized for acute rectal loading or impaction.
    \item \textbf{Medication review:} Reducing or substituting drugs that cause constipation where clinically feasible.
\end{itemize}

\section*{Complications}
\begin{itemize}
    \item \textbf{Electrolyte imbalances:} Chronic overuse can lead to hypokalaemia, hyponatraemia, and metabolic acidosis or alkalosis.
    \item \textbf{Dehydration:} Excessive loss of water in watery stools can lead to volume depletion and renal impairment.
    \item \textbf{Laxative dependency:} The colon may become less responsive to normal physiological stimuli (melanosis coli or ``cathartic colon''), making natural bowel movements difficult.
    \item \textbf{Malabsorption:} Rapid intestinal transit can impair the absorption of nutrients, vitamins, and oral medications.
    \item \textbf{Haemorrhoids and anal fissures:} Persistent straining before or during inadequate laxative action can cause mucosal tears and vascular swelling.
\end{itemize}

\section*{Prognosis and Outcomes}
The prognosis for patients requiring short-term laxatives is excellent, with most experiencing rapid relief and returning to normal bowel function once lifestyle factors are optimised or the offending secondary cause is removed. Chronic laxative users or those with severe underlying motility disorders may require ongoing, multidisciplinary management and long-term care to prevent dependency and complications. Overall, correct utilisation of laxatives under medical supervision leads to improved quality of life without long-term adverse sequelae.

\section*{Prevention and Self-Care}
\begin{itemize}
    \item \textbf{Gradual dietary changes:} Introduce high-fibre foods slowly to prevent bloating and flatulence.
    \item \textbf{Hydration:} Drink plenty of water throughout the day, especially when using bulk-forming laxatives.
    \item \textbf{Regular toilet habits:} Respond promptly to the urge to defaecate, particularly after meals when the gastrocolic reflex is active.
    \item \textbf{Physical activity:} Aim for at least 30 minutes of moderate exercise daily to stimulate intestinal motility.
    \item \textbf{Avoid chronic stimulant use:} Use stimulant laxatives only for short-term relief (less than a week) unless prescribed by a doctor.
\end{itemize}

\section*{Sources and Further Reading}
The following reputable organisations provide further information on laxatives and bowel health:
\begin{itemize}
    \item Healthdirect Australia. \textit{Laxatives overview and safety guidelines.} https://www.healthdirect.gov.au/
    \item Better Health Channel. \textit{Constipation and laxative use information.} https://www.betterhealth.vic.gov.au/
    \item Mayo Clinic. \textit{Over-the-counter laxatives for constipation: Use with caution.} https://www.mayoclinic.org/
    \item NHS. \textit{Laxatives usage, types, and side effects.} https://www.nhs.uk/
    \item MSD Manual. \textit{Laxatives classification and clinical pharmacology.} https://www.msdmanuals.com/
\end{itemize}
